\begin{abstract}
Neural networks can continue learning while successive gradients nearly reverse direction. Understanding when this oscillation remains productive requires connecting local optimization geometry to accumulated parameter change and to movement of the learned function. We characterize constructive and destructive training regimes through gradient alignment and magnitude, distinguishing local reversal from cumulative displacement and loss reduction. Under smooth full-batch gradient descent, a scalar computed from two existing gradients equals the step-averaged directional Hessian curvature; we give its pointwise approximation error and explicit sampling-noise contribution. Controlled experiments with multilayer perceptrons, residual networks, and Transformers show that negative adjacent interaction can coexist with positive total cross-step interaction and lower training loss. We then interpret the same observable under Adam as a projected gradient response and evaluate its association with future loss increases, retaining missed events and false alarms. A fixed-probe tracing algorithm reveals concentrated representation motion, shared temporal structure across depth, and repeated reversals along leading output patterns. Finally, a common-state comparison of learning-rate reduction, second-moment changes, gradient clipping, and local affine-displacement caps tests the functional interpretation. Across three single-initialization cases, the preselected median cap reduces mean training loss by 32.2\%, 9.4\%, and 19.8\%, while reversals persist. Endpoint and held-out tradeoffs limit the conclusion to controlled training dynamics. The resulting framework connects interaction, curvature, and functional movement without equating oscillation with failed learning.
\end{abstract}
